\section{Video-coding Principles}
We want to exploit both spatial redundancy and temporal redundancy.
\subsection{Block Schema}
$$x \longrightarrow \boxed{\text{Temporal Compression}} \longrightarrow \boxed{\text{Spatial Compression}} \longrightarrow \boxed{\text{Buffer}} \longrightarrow$$
Can we swap Temporal and Spatial compression blocks? Yes but it isn't used because it doesn't make sense.
\begin{center}
\includegraphics[width=0.8\linewidth]{hybrid_concept.png}
\end{center}
\noindent Hybrid coding intuition: only the temporal prediction error $e = x - x_p$ (previous frame) is quantized and transmitted; the reconstructed frame $x_q$ is stored in the frame buffer to predict the next one (closed loop, same anti-drift principle as DPCM).
\subsection{Block-Matching motion estimation}
$$J(v) = d(B_k^{(p)},B_h^{(p+v)}) + \lambda_{ME}R(v)$$
How can we build the prediction using motion estimation? Motion Compensation
\subsubsection{Motion Compensation}
$$\hat{I}_k(p) = I_h(p + v^\star(p))$$
Sometimes MC works better and sometimes it works worse, solution: adaptive blocks.
\subsubsection{Adaptive block coding}
\begin{itemize}
    \item Intra Coding: encode block as it is
    \item Inter Coding: use other images to predict it
\end{itemize}
$\forall B_k^{(p)}$ block of image k:
\begin{itemize}
    \item compute $J(v)$
    \item find the $v^\star = \arg\min_pv(p)$
    \item encode block:
    \begin{itemize}
        \item If intra: encode in a JPEG-like way
        \item If Inter:
        \begin{itemize}
            \item Decode $v^\star$ motion vector
            \item Decode predictor error blocks $E_p^{(b)}$
            \item Correct blocks
        \end{itemize}
    \end{itemize}
\end{itemize}
\subsubsection{Design Parameters}
Block size (smaller better), Cost function (SAD + regularization),\\
Motion Model (translational), Search Strategy (fast methods: hex, TZ)
\subsection{GOP: Group of Pictures}
It is a sequence of $n$ frames, usually in the following order:
$$\boxed{\boxed{I}\ \boxed{B}\ \boxed{B}\ \boxed{P}\ \boxed{B}\ \boxed{B}\ \boxed{P}\ \boxed{B}\ \boxed{B}} \to n\text{-th GOP}$$
\subsubsection{`I' frames: Intra Frames}
Access points of our videos, decoding starts here. Useful to stop error propagation. Fast to encode, no ME.
% Integrated from Summaries
All blocks of an I frame are Intra coded. Used for: \textbf{random access} (independently decodable), \textbf{fast forward} (decode only I frames), \textbf{error robustness} (stop error propagation). Low complexity but low compression.
\subsubsection{`B' frames: `between' frames}
Uses both past and future frame exploitation for prediction. This leads a change in encoding order (1,4,2,3,7,5,6,10,8,9), leading also to a \textbf{structural latency}.
In Real time applications we do not use B frames due to this latency.
% Integrated from Summaries
For each block of a B frame the encoder can choose: Intra coding, \textbf{forward} prediction (from the past), \textbf{backward} prediction (from the future), or \textbf{bidirectional} prediction (average of past and future predictors). Very high complexity (double ME) but the highest compression ratio.
\subsubsection{`P' frames: Predictive Frames}
Predicted from the \textbf{previous anchor frame (I or P)}. Decision taken at block level: each block can be Inter or Intra coded. High complexity (ME/MC) but much higher compression than I frames at the same quality.
\subsubsection{Rate and quality of the frame types}
I frames are $\approx 3\to5\times$ larger than P frames and $\approx 10\to20\times$ larger than B frames. The quality of I frames is typically forced high, because they predict (directly or indirectly) the whole GOP.\\
The GOP structure controls the tradeoff among: compression efficiency, delay, random access, error propagation.
\subsection{Hybrid Video Encoder}
\subsubsection{Frame coding modes}
Encoder tries all of these to encode a frame:
\begin{itemize}
    \item Intra: available for all frames
    \item Inter: ME/MC-based temporal prediction (only NON-intra frames)
    \item Direct (skip): trying to infer MV from neighbors (gain in rate, loss in distorsion)
    \item Lossless: all frames can be encoded in this way
\end{itemize}
\subsubsection{Block Size: Block Partition Problem}
$$D = \sr{K}{k=1}D_k(i_k, Q) \qquad \qquad R = \sr{K}{k=1}R_k(i_k, Q)$$
We want the lowest distorsion for a given bitrate:
$$J(\vec{i},Q,\lambda ) = D(\vec{i},Q) + \lambda R(\vec{i},Q) \qquad \text{ where }\vec{i} = \left\{i_k\right\}_{k=1}^K \text{ s.t. }i_k^\star = \arg\min_{i_k} J(i_k,Q,\lambda) = \arg\min_{i_k} J_k(\vec{i})$$
% Integrated from Summaries
How to find Q and $\lambda$? $Q$ is an input (set by rate control); for each $Q$ there is an optimal $\lambda$, determined \textbf{empirically} per codec:
$$\text{MPEG-2: } \lambda = aQ^2 + b \qquad \text{H.264: } \lambda = c\cdot2^{dQ+e} \qquad \lambda_{ME} = \sqrt{\lambda}$$
Joint minimization over all $\vec{i}$ is too complex $\Rightarrow$ suboptimal \textbf{block-wise} minimization of $J_k(i_k,Q,\lambda) = D_k + \lambda R_k$, one block at a time. Block partition cast as a mode too: split if $J_{split} = \sr{}{i}J_{subblock_i} < J_B$, applied recursively from the largest block size.
\subsubsection{Coding Mode $i_k$ selection}
On a R-D plane we plot all the points related to the available modes, then we select the nearest to our J slope.
% Integrated from Summaries
$D = -\lambda R + J$ is a family of lines with slope $-\lambda$: high $\lambda$ (steep) $\to$ low bit budget, picks low-rate modes (e.g. Direct); low $\lambda$ (flat) $\to$ prioritizes fidelity (e.g. Inter8 bidirectional). The chosen mode is the first point touched by the line moving from the origin towards the convex hull of the available modes.
% Integrated from Summaries: scheme to be able to draw at the exam
\subsubsection{Encoder scheme}

\begin{center}
\includegraphics[width=\linewidth]{hybrid_encoder.png}
\end{center}
\noindent Full hybrid encoder block diagram. The dashed JPEG-like core (Transform $\to$ Quantization $\to$ Lossless Coding) compresses the residual $e_k = B_k - \tilde{B}_k$. The reconstruction loop (Inv.\ Quantization $\to$ Inverse Transform $\to$ $+\,\tilde{B}_k$) feeds the Frame Buffer used by Motion Compensation, Motion Estimation and Intra Prediction; Mode Decision picks the predictor and Control drives rate via the Channel Buffer.
Comment on the blocks:
\begin{itemize}
    \item Each image is split into blocks; for each block the encoder computes a prediction $v$: \textbf{null} (direct DCT coding), \textbf{temporal} (ME/MC, simple or bi-prediction), or \textbf{spatial} (Intra, from already encoded blocks of the same frame).
    \item The prediction error $y = x_k - v$ is JPEG-like compressed: Transform + Quantization.
    \item \textbf{Decoder loop inside the encoder}: $Q^{-1}$, $T^{-1}$, then $\tilde{x}_k = \hat{y} + v$ is stored in the \textbf{DFB} (Decoded Frame Buffer), because future temporal/spatial predictions must use the same reconstructed data the decoder will have (no drift, same principle as closed-loop DPCM).
    \item Side information (mode decision, block partition, motion vectors) + quantized coefficients are entropy-coded with a VLC (Huffman/Arithmetic/Exp-Golomb; MVs coded predictively, exploiting spatial coherence of the MVF).
    \item \textbf{Rate control}: encoded stream enters the channel buffer at rate $R_C$ (content-dependent) and leaves at target rate $R_T$. Closed-loop controller: if occupancy $> \gamma_{high}$ increase quantization step ($R_C$ drops); if $< \gamma_{low}$ decrease it ($R_C$ grows).
\end{itemize}
Why do we have standards? Interoperability, we only standardize the decoder (able to decode anything), to leave encoding competition.
\subsection{Hybrid Video Decoder}
\subsubsection{Decoder scheme}
% Integrated from Summaries
\begin{center}
\includegraphics[width=0.85\linewidth]{hybrid_decoder.png}
\end{center}
\noindent Hybrid decoder: only the reconstruction half of the encoder. After Lossless Decoding, Inv.\ Quantization and Inverse Transform give $\tilde{e}_k$, summed with the prediction (Motion Compensation from the Frame Buffer, or Intra Prediction, selected by Mode) to output $\tilde{B}_k$. No mode decision / motion estimation $\Rightarrow$ far cheaper than the encoder.
\begin{itemize}
    \item First decodes the \textbf{side information}: partition, coding mode, motion vector(s).
    \item Then inverse quantization, inverse transform, and sum with the prediction (MC or Intra) to get the final block, stored back in the DFB.
    \item The decoder does \textbf{not} perform: mode decision, block partition decision, motion estimation (the heavy operations) $\Rightarrow$ decoder much faster than encoder. These are read from the bitstream.
    \item Since only the decoder is standardized, ME strategy / mode decision / partitioning are free: competition on encoder efficiency (hex-search, fast mode decision, \dots).
\end{itemize}
\subsubsection{Key elements}
\begin{itemize}
    \item Spatial Compression: Transforms
    \item Temporal Compression: GOP and ME
    \item LossLess coding: Huffman and Arithmetic codings
    \item Coding Mode estimation: Lagrangian ($D + \lambda R$)
\end{itemize}
\subsection{Video Encoding Standards}
\begin{itemize}
    \item MPEG (1,2,4): among the first standards, since 1988
    \item H.264/AVC: free standard
    \item H.265/HEVC: some patenting
    \item H.266/VVC: not really free, we have some patenting
    \item AV1, VP8, VP9: Royalty free standards (AOM\footnote{Alliance for Open Media})
\end{itemize}
\subsubsection{MPEG-1}
MPEG-1 Part 2 (MP3 Audio Format) was the first digital audio standard used for CDs.\\
MPEG-1 allows:
\begin{itemize}
    \item Rate $\leq 1.86$Mbps
    \item $[720\times 576]@30$fps max resolution
    \item Audio: 128 $\to$ 320 kbps
\end{itemize}
